%% source: 2021-fa-midterm_01
%% tags: [binary search trees]
\begin{prob}
    Define the \textbf{smallest gap} in a collection of numbers to be the smallest difference
    between two distinct elements in the collection (in absolute value).
    For example, the smallest gap in $\{4, 9, 1, 6\}$ is 2 (between 4 and 6).

    Suppose a collection of $n$ numbers is stored in a \textbf{balanced} binary search
    tree. 
    What is the time complexity required for an efficient algorithm to calculate
    the smallest gap of the numbers in the BST? State your answer as a function of $n$ in
    asymptotic notation.

    \begin{soln}
        $\Theta(n)$. The smallest gap in the BST can occur between any two consecutive pairs of numbers 
        in the BST, meaning that no matter what we need to search through the whole tree to obtain this 
        number, taking $\Theta(n)$ time.
    \end{soln}

\end{prob}
